\chapter{Implementation}
\label{chap:imp}
\lhead{\emph{Project Implementation}}

\section{Introduction}

This chapter summarizes the implementation work completed for the supervised and semi-supervised learning benchmark framework. The focus is on practical delivery: what was built, what problems were encountered, and how those problems were resolved in the final project codebase.

\section{What Is Implemented}

The current project implementation includes the following runnable training components:

\begin{itemize}
    \item Supervised baseline training
    \item Pseudo-Label training
    \item FixMatch training
    \item MixMatch training
    \item FlexMatch training
\end{itemize}

These are implemented through dedicated trainer modules under \texttt{src/training/} and executed through the project training scripts (for example, \texttt{train\_supervised.py}, \texttt{train\_supervised\_limited.py}, and \texttt{train\_ssl.py}).

The implementation also includes:

\begin{itemize}
    \item YAML-based configuration files for supervised and SSL runs
    \item Label-budget protocols for CIFAR-10 (250, 1000, 4000 labels per class)
    \item Checkpoint save/resume support (best and last checkpoints)
    \item Log generation for reproducibility and debugging
    \item Docker-based execution support with GPU passthrough
\end{itemize}

\section{Difficulties Encountered and Resolutions}

This section keeps the selected key difficulties and documents the practical resolution for each.

\subsection{Difficulty 1: CUDA Infrastructure and Docker GPU Passthrough}

\textbf{Classification:} Hard \quad \textbf{Timeline Consumed:} 1 week

\textbf{Problem:}
Docker GPU setup required correct NVIDIA runtime/toolkit configuration and explicit GPU flags. Missing or inconsistent setup caused unintended CPU training and significant slowdown.

\textbf{Resolution Implemented:}

\begin{itemize}
    \item Added pre-run GPU verification using \texttt{verify\_installation.py}
    \item Added Docker helper scripts that enforce GPU flags
    \item Added startup logging to confirm runtime device and hardware details
\end{itemize}

\subsection{Difficulty 2: Stop-Loss Threshold and Warmup Interaction}

\textbf{Classification:} Medium \quad \textbf{Timeline Consumed:} 3 days

\textbf{Problem:}
Stop-loss checks could terminate runs too early in low-label regimes due to unstable validation loss.

\textbf{Resolution Implemented:}

\begin{itemize}
    \item Introduced warmup-aware stop-loss checks
    \item Used conservative stop-loss behavior for low-label runs
    \item Added explicit trigger logging for diagnosis
    \item Standardized settings through configuration templates
\end{itemize}

\subsection{Difficulty 3: Per-Class Accuracy Variance and Data Stratification}

\textbf{Classification:} Medium \quad \textbf{Timeline Consumed:} 3--5 days

\textbf{Problem:}
Random label selection created class imbalance at low-label budgets, reducing reliability of per-class comparisons.

\textbf{Resolution Implemented:}

\begin{itemize}
    \item Enforced per-class label budgets in data splitting
    \item Added stratified validation handling
    \item Logged class distributions for each run
\end{itemize}

\subsection{Difficulty 5: Time Constraints and Local Hardware Limits}

\textbf{Classification:} Hard \quad \textbf{Timeline Consumed:} Ongoing

\textbf{Problem:}
Full experiment coverage (algorithms, budgets, and multi-seed repetition) exceeded available local compute time.

\textbf{Resolution Implemented:}

\begin{itemize}
    \item Prioritized core algorithms first, then extended to FlexMatch
    \item Used single-seed initial runs for implementation validation
    \item Used checkpoint/resume and queued overnight execution
    \item Reduced scope to stable, fully implemented methods
\end{itemize}

\subsection{Difficulty 7: MixMatch Instability at 250 Labels}

\textbf{Classification:} Medium \quad \textbf{Timeline Consumed:} 3--4 days

\textbf{Problem:}
MixMatch showed unstable behavior in the 250-label setting under fixed project configuration.

\textbf{Resolution Implemented:}

\begin{itemize}
    \item Treated MixMatch-250 as a limitation case in this project
    \item Used validation monitoring and early stopping safeguards
    \item Prioritized more stable methods for low-label conclusions
\end{itemize}

\subsection{Difficulty 8: Multi-Seed Reproducibility Constraints}

\textbf{Classification:} Medium \quad \textbf{Timeline Consumed:} Deferred

\textbf{Problem:}
Running full multi-seed experiments for all configurations was computationally infeasible on a single local GPU within timeline constraints.

\textbf{Resolution Implemented:}

\begin{itemize}
    \item Kept seed as an explicit configurable parameter
    \item Scheduled selective multi-seed runs for highest-priority comparisons
    \item Deferred complete multi-seed coverage as future work
\end{itemize}

\subsection{Difficulty 9: Data Augmentation Consistency Across Algorithms}

\textbf{Classification:} Easy \quad \textbf{Timeline Consumed:} 1 day

\textbf{Problem:}
Hardcoded augmentation logic made algorithm-specific behavior difficult to maintain and compare consistently.

\textbf{Resolution Implemented:}

\begin{itemize}
    \item Moved augmentation control into configurable components
    \item Enabled per-algorithm weak/strong augmentation handling
    \item Improved consistency of augmentation behavior across runs
\end{itemize}

\subsection{Difficulty 10: Checkpoint Management and File I/O Robustness}

\textbf{Classification:} Medium \quad \textbf{Timeline Consumed:} 2--3 days

\textbf{Problem:}
Interrupted runs and file-write issues could produce invalid checkpoints and failed resume attempts.

\textbf{Resolution Implemented:}

\begin{itemize}
    \item Added safer checkpoint writing and loading paths
    \item Kept both \texttt{best} and \texttt{last} checkpoints
    \item Added clearer error handling for resume failures
\end{itemize}

\section{Actual Solution Approach vs. Original Design}

This section contrasts the original plan with what was actually delivered in the final implementation.

\subsection{Algorithms: Planned vs. Implemented}

\subsubsection{Original Plan}

The design targeted a broad SSL comparison with core methods and additional methods if time permitted.

\subsubsection{Actual Implementation}

The final project implements the following methods in one shared framework:

\begin{itemize}
    \item Supervised baseline
    \item Pseudo-Label
    \item FixMatch
    \item MixMatch
    \item FlexMatch
\end{itemize}

Deferred methods: Temporal Ensembling, Mean Teacher, FreeMatch, SoftMatch, SimMatch, and UDA/VAT variants.

\subsubsection{Justification}

\begin{enumerate}
    \item Infrastructure and stability issues (D1, D2, D7) reduced available implementation time.
    \item Compute constraints (D5, D8) made full 8--10 method coverage impractical.
    \item A correctness-first strategy was chosen: fewer fully working methods were prioritized over partially implemented methods.
\end{enumerate}

\subsection{Datasets: Planned vs. Actual}

\subsubsection{Original Plan}

Primary focus was CIFAR-10, with additional datasets considered if time allowed.

\subsubsection{Actual Implementation}

The framework remains extensible, but the completed benchmark work in this project focuses on CIFAR-10.

\begin{enumerate}
    \item CIFAR-10 provides a standard SSL benchmark for fair method comparison.
    \item Additional datasets would substantially increase total training time.
    \item The selected label budgets (250, 1000, 4000 per class) were sufficient to evaluate low, medium, and high label regimes.
\end{enumerate}

\subsection{Evaluation Setup: Planned vs. Actual}

\subsubsection{Original Plan}

The plan considered a wider set of label fractions with stronger statistical replication.

\subsubsection{Actual Implementation}

Evaluation was standardized around three CIFAR-10 label budgets:

\begin{itemize}
    \item 250 labels per class
    \item 1000 labels per class
    \item 4000 labels per class
\end{itemize}

This setup balanced methodological usefulness with available compute resources.

\subsection{Infrastructure: Planned vs. Actual}

\subsubsection{Original Plan}

The design allowed for portable container-based execution, including possible cloud scaling.

\subsubsection{Actual Implementation}

The final workflow is optimized for reproducible single-machine execution with Docker and local GPU usage.

\begin{itemize}
    \item Verification-first runtime checks were added (D1).
    \item Sequential and overnight execution was used for throughput under local constraints (D5).
    \item Checkpoint-resume logic improved recovery from interruptions (D10).
\end{itemize}

\subsection{Evaluation Outputs Implemented}

The final codebase produces implementation-ready outputs for analysis and reporting:

\begin{itemize}
    \item Per-run logs and checkpoint artifacts
    \item Per-run accuracy/loss histories
    \item Comparison summaries in JSON/CSV/Markdown outputs
    \item Calibration and confusion diagnostics for supervised vs SSL comparison runs
\end{itemize}

Full multi-seed statistical reporting remains partial and is documented as future work (D8).

\subsection{Original Risk Assessment: Outcomes}

The original risk categories were largely accurate in practice:

\begin{itemize}
    \item Risk 2 (limited resources) materialized and drove scope reduction (D5, D8)
    \item Risk 3 (SSL complexity) materialized through instability and augmentation consistency issues (D7, D9)
    \item Risk 5 (environment/integration) materialized through Docker/CUDA setup complexity (D1)
\end{itemize}

\section{Summary and Conclusion}

The delivered implementation is a focused, working SSL benchmark framework with five integrated methods and a reproducible training workflow. The selected difficulties were resolved through practical engineering decisions: environment verification, warmup-aware training safeguards, stratified data handling, scope control under compute limits, and robust checkpoint management.

The project therefore meets its implementation objective: a stable comparative framework that is ready for evaluation in the current scope and extensible for future algorithm additions.
